\documentclass[12pt,a4paper]{article}

% ---------------------------------------------------------------
% Packages
% ---------------------------------------------------------------
\usepackage[utf8]{inputenc}
\usepackage[T1]{fontenc}
\usepackage[english]{babel}
\usepackage{lmodern}
\usepackage{microtype}

\usepackage{amsmath,amssymb,amsthm}
\usepackage{mathtools}
\usepackage{bm}

\usepackage{geometry}
\geometry{margin=2.5cm, top=3cm, bottom=3cm}

\usepackage{booktabs}
\usepackage{array}
\usepackage{multirow}
\usepackage{tabularx}

\usepackage{hyperref}
\hypersetup{
  colorlinks=true,
  linkcolor=black,
  citecolor=black,
  urlcolor=blue,
  pdftitle={Extended Security Analysis of the Alaniz Cipher},
  pdfauthor={Alfredo Rodriguez Langa},
}

\usepackage{xcolor}
\usepackage{fancyhdr}
\usepackage{titlesec}
\usepackage{abstract}

\usepackage{listings}
\lstset{basicstyle=\ttfamily\small, frame=single, breaklines=true}

% ---------------------------------------------------------------
% Theorem environments
% ---------------------------------------------------------------
\newtheorem{theorem}{Theorem}[section]
\newtheorem{lemma}[theorem]{Lemma}
\newtheorem{corollary}[theorem]{Corollary}
\newtheorem{proposition}[theorem]{Proposition}
\theoremstyle{definition}
\newtheorem{definition}[theorem]{Definition}
\newtheorem{remark}[theorem]{Remark}

% ---------------------------------------------------------------
% Math macros
% ---------------------------------------------------------------
\newcommand{\Fp}{\mathbb{F}_p}
\newcommand{\Fq}{\mathbb{F}_q}
\newcommand{\GL}{\mathrm{GL}}
\newcommand{\Hzero}{H^{0}}
\newcommand{\Img}{\mathrm{Im}}
\newcommand{\Tr}{\mathrm{Tr}}
\newcommand{\bsigma}{\bm{\sigma}}
\newcommand{\calF}{\mathcal{F}}
\newcommand{\norm}[1]{\lVert #1 \rVert}

% ---------------------------------------------------------------
% Header / footer
% ---------------------------------------------------------------
\pagestyle{fancy}
\fancyhf{}
\rhead{\small Extended Security Analysis of the Alaniz Cipher}
\lhead{\small A.~Rodr\'iguez Langa, INECO}
\cfoot{\thepage}
\renewcommand{\headrulewidth}{0.4pt}

% ---------------------------------------------------------------
% Title
% ---------------------------------------------------------------
\title{%
  \textbf{Extended Security Analysis of the Alaniz Cipher}\\[6pt]
  \large Chosen-Plaintext Attacks, Key Recovery, and Structural\\
  Properties of the Local Encryption Map\\[4pt]
  \normalsize\textit{Companion analysis to: L.~Alaniz Pintos,
  ``The Alaniz Cipher: Encryption from Nonlinear Sheaf Morphisms
  over Graphs'', 2026}
}

\author{%
  Alfredo Rodr\'iguez Langa\\[4pt]
  \normalsize Smart Products Division, INECO\\
  \normalsize Madrid, Spain\\[2pt]
  \small \href{mailto:alfredo.rodriguez@ineco.com}{alfredo.rodriguez@ineco.com}
}

\date{March 2026}

% ---------------------------------------------------------------
% Document
% ---------------------------------------------------------------
\begin{document}

\maketitle
\thispagestyle{fancy}

% ---------------------------------------------------------------
\begin{abstract}
This document reports seven reproducible security experiments
(Experiments~07--13) that complement and stress-test the security
claims of the Alaniz Cipher (SHEC). The experiments probe the local
encryption map $g_B(y) = y + B\,\sigma(y)$ for collision properties,
attempt key recovery under known- and chosen-plaintext models, and
analyse the algebraic structure exploitable by an adversary.

The central finding is contained in \textbf{Experiment~10}: both
supported nonlinear maps ($\sigma = x^3$ and $\sigma = x^{p-2}$)
satisfy a componentwise scaling law
$\sigma(\lambda y)_i = \lambda^e\,\sigma(y)_i$, which allows a
chosen-plaintext attacker to recover the per-node key matrices
$(A_v, B_v)$ using only $O(d)$ encryption queries---independently of
the field size~$p$ and the size of $\GL(d,\Fp)$.
This result directly contradicts the CPA-resistance claim in the
original paper and is not addressed by the NL-SMIP hardness
reduction, which concerns a distinct adversarial model.
The remaining experiments characterise the collision structure of
$g_B$ (Experiments~07 and~11), measure the concrete complexity of
exhaustive and differential key recovery under the known-plaintext
model (Experiments~08, 09, 12, and~13), and provide the algebraic
ground-truth necessary for evaluating any proposed fix.
\end{abstract}

\tableofcontents
\newpage

% ================================================================
\section{Introduction}
% ================================================================

The Alaniz Cipher encrypts a message encoded as a global section
$s \in \Hzero(G, \calF)$ of a cellular sheaf over a graph~$G$.
At each node $v$ the local ciphertext is
\begin{equation}\label{eq:encrypt}
  c_v = A_v \cdot s_v + B_v \cdot \sigma(A_v \cdot s_v),
  \qquad (A_v, B_v) \in \GL(d, \Fp)^2,
\end{equation}
where $\sigma : \Fp^d \to \Fp^d$ is a public componentwise nonlinear
map. The original paper argues that this construction resists
chosen-plaintext attacks (CPA) because $\sigma$ defeats linearisation:
any attempt to model $c_v$ as a linear function of $s_v$ produces an
inconsistent system as soon as enough ciphertext pairs are collected.

The experiments in this document examine whether non-linear attacks
succeed. They are organised in order of increasing algebraic
sophistication:

\begin{itemize}
  \item \textbf{Experiments~07 and~11} characterise the map
    $g_B(y) = y + B\,\sigma(y)$, the core function the attacker
    must invert.
  \item \textbf{Experiment~08} establishes the exact threshold at
    which the key becomes trivially recoverable ($d = 1$) and
    confirms algebraically that linear attacks fail for $d \geq 2$.
  \item \textbf{Experiments~09 and~12} measure the empirical cost of
    key recovery under the known-plaintext model via exhaustive
    A-search.
  \item \textbf{Experiment~10} presents a direct CPA break requiring
    only $O(d)$ queries.
  \item \textbf{Experiment~13} proposes a differential column attack
    that recovers the columns of $A$ independently, reducing the
    search space from $|\GL(d,p)|$ to $d\,(p^d - 1)$ candidates.
\end{itemize}

All experiments use the production implementation in \texttt{alaniz/}
without modification. Random seeds are fixed for full reproducibility.

% ================================================================
\section{Protocol Notation}
% ================================================================

For reference, Table~\ref{tab:notation} summarises the quantities
used throughout this document.

\begin{table}[h]
\centering
\caption{Notation used in this analysis.}
\label{tab:notation}
\begin{tabularx}{\linewidth}{>{\ttfamily}lX}
\toprule
\normalfont Symbol & Meaning \\
\midrule
$p$                & Field characteristic (prime) \\
$d$                & Fibre dimension \\
$G$                & Graph with $n$ nodes \\
$\Hzero(G, \calF)$ & Space of global sections (plaintext space) \\
$s_v \in \Fp^d$    & Local plaintext at node $v$ \\
$(A_v, B_v)$       & Secret key at node $v$; both in $\GL(d,\Fp)$ \\
$\sigma$           & Public componentwise map:
                     \texttt{inverse} ($x_i \mapsto x_i^{p-2}$)
                     or \texttt{cube} ($x_i \mapsto x_i^3$) \\
$y_v = A_v s_v$    & Secret intermediate vector at node $v$ \\
$g_B(y)$           & $y + B\,\sigma(y)$: local encryption function \\
\bottomrule
\end{tabularx}
\end{table}

Decryption solves $y + B_r\,\sigma(y) = c_r$ at the root node,
propagates via restriction maps, and uses cohomological filtering to
select the valid global section.

% ================================================================
\section{Experiment 07 --- Collision Table for the Local Encryption Map}
% ================================================================

\subsection{Objective}

Quantify the surjectivity and preimage multiplicity of $g_B$ over the
full domain $\Fp^d$ for $d \in \{1, 2, 4\}$ and
$p \in \{17, 23, 29, 47, 59, 89, 131, 251\}$.
This provides concrete evidence of whether $g_B$ is injective, which
determines decryption ambiguity.

\subsection{Methodology}

For $d \in \{1, 2\}$, all $p^d$ inputs are enumerated exactly.
The preimage map $\{c \mapsto g_B^{-1}(c)\}$ is constructed and the
following metrics are computed per random $B \in \GL(d, \Fp)$
(averaged over eight independent draws):
\begin{equation}
  \text{image fraction} = \frac{|\Img(g_B)|}{p^d}, \qquad
  \gamma_B = \sum_{c \in \Fp^d} \!\left(\frac{|g_B^{-1}(c)|}{p^d}\right)^{\!2}.
\end{equation}
For $d = 4$, where exhaustive enumeration is infeasible, a birthday
probe with $3\lceil\sqrt{p^d}\rceil$ random samples guarantees a
collision with high probability by the birthday bound. Every collision
example is verified against the algebraic condition
$\Delta y = B\bigl(\sigma(y_2) - \sigma(y_1)\bigr)$.

\subsection{Findings}

Across all tested primes and both $\sigma$ functions, the image
fraction stabilises near $0.53$ for $\sigma = \texttt{inverse}$
and $0.66$ for $\sigma = \texttt{cube}$ in dimension $d = 1$, and
near $0.62$--$0.64$ for $d = 2$. The map $g_B$ is therefore
consistently non-surjective: roughly one third of ciphertexts at
each node are unreachable under a fixed key. Collisions are confirmed
at all parameter sizes, including $d = 4$, $p = 251$.

\subsection{Security Implications}

Non-surjectivity implies two consequences. First, decryption at the
root node admits multiple preimages (at most the maximum multiplicity,
observed at 2--5 for $d = 2$), introducing ambiguity that the
cohomological filter resolves. Second, the collision probability
$\gamma_B$ provides a direct empirical estimate of the parameter
$\gamma$ from Theorem~7.1 of the original paper.

% ================================================================
\section{Experiment 08 --- Key Recovery at $d=1$ and Linear Attack Diagnostics}
% ================================================================

\subsection{Objective}

Establish the exact threshold below which the key is algebraically
recoverable, and confirm that the linearisation defence holds for
$d \geq 2$.

\subsection{Exact Recovery for $d = 1$}

For $d = 1$ the encryption reduces to a scalar equation. Given two
plaintext--ciphertext pairs $(s_1, c_1)$ and $(s_2, c_2)$:

\paragraph{Case $\sigma = \texttt{inverse}$.}
Multiplying $c_i = a s_i + b(as_i)^{-1}$ through by $as_i$ and
subtracting the two equations yields:
\begin{equation}
  a = \frac{c_1 s_1 - c_2 s_2}{s_1^2 - s_2^2}, \qquad
  b = c_1(as_1) - (as_1)^2.
\end{equation}

\paragraph{Case $\sigma = \texttt{cube}$.}
The two equations form a $2\times 2$ Vandermonde-type system in
$(a,\, ba^3)$, solved by Cramer's rule with determinant
$s_1 s_2^3 - s_2 s_1^3$.

Both formulas are exact over $\Fp$ and require only two captures.

\subsection{Linear Surrogate Failure for $d \geq 2$}

For $d \in \{2, 4\}$, the experiment builds the linear surrogate
$c_v = M s_v$ (treating the key as a single $d\times d$ matrix $M$)
and tracks its rank and consistency as captures accumulate. As
expected, the nonlinearity of $\sigma$ renders the system
inconsistent after enough pairs, confirming that pure linearisation
fails.

\subsection{Findings and Security Implications}

Dimension $d = 1$ provides \textbf{no computational security}: the
key is recoverable with 2 plaintext-ciphertext pairs using only
modular arithmetic. The minimum secure dimension is $d \geq 2$.
However, Experiment~10 shows that a non-linear CPA attack succeeds
at all dimensions.

% ================================================================
\section{Experiment 09 --- Exhaustive Known-Plaintext Key Recovery}
% ================================================================

\subsection{Objective}

Measure the concrete cost of recovering the per-node key $(A_v, B_v)$
under the known-plaintext model via exhaustive search over
$A \in \GL(d, \Fp)$ with linear solve for $B$.

\subsection{Methodology}

For each candidate matrix $\hat{A}$, compute
$y_i = \hat{A}\, s_i$ and $z_i = \sigma(y_i)$ for $d$ fit captures,
then solve the linear system
\begin{equation}
  c_i - y_i = B\, z_i, \qquad i = 1, \ldots, d,
\end{equation}
for the $d^2$ entries of $B$. Validate the recovered $(\hat{A}, \hat{B})$
against all $k$ captures. The campaign covers
$p \in \{17, 23, 29, 47, 59, 89, 131, 251\}$, $d = 2$,
with $k \in \{4, 8, 12\}$ captures. For
$|\GL(2, p)| \leq 10^5$ exhaustive enumeration is used;
for larger primes 5\,000 random samples are drawn.

\subsection{Findings and Security Implications}

For $p \leq 29$, $|\GL(2, p)| \leq 7\times 10^5$, and the true key
is recovered exactly in all trials. For $p \geq 47$ the sample
probability is negligible (${\approx}5000/|\GL(2,p)|$). The number
of false positives drops sharply with $k$: at $k = 12$ no false
positive survives any exhaustive run. The exhaustive A-search is
intractable at standard parameters ($d = 4$, $p \approx 2^{31}$),
but Experiment~10 renders this bound irrelevant by recovering the
key in $O(d)$ chosen-plaintext queries regardless of $p$.

% ================================================================
\section{Experiment 10 --- Direct CPA Key Recovery via Scaling Homogeneity}
% ================================================================

\textit{This is the main security finding of this analysis.}

\subsection{Objective}

Demonstrate a chosen-plaintext attack that recovers $(A_v, B_v)$
in $O(d)$ encryption queries, without any exhaustive search over
$\GL(d, \Fp)$.

\subsection{Mathematical Basis}
\label{sub:scaling}

Both supported nonlinear maps satisfy a \textbf{componentwise scaling
law}:
\begin{equation}\label{eq:scaling}
  \sigma(\lambda\, y)_i = \lambda^e\, \sigma(y)_i
  \qquad \forall\, y \in \Fp^d,\; \lambda \in \Fp^*,
\end{equation}
with $e = 3$ for $\sigma = \texttt{cube}$ and $e \equiv p-2 \pmod{p-1}$
for $\sigma = \texttt{inverse}$.

Because $B$ acts linearly on $\sigma(y)$, equation~\eqref{eq:scaling}
extends to the full ciphertext. For plaintexts $s = t\,e_j$
(scalar multiples of the $j$-th canonical basis vector of the
$H^0$-coefficient space), the local ciphertext satisfies
\begin{equation}\label{eq:pair}
  c(t\,e_j) = t\,a_j + t^e\,B\,\sigma(a_j),
\end{equation}
where $a_j = A\,e_j$ is the $j$-th column of $A$.
Querying $s_1 = e_j$ ($t = 1$) and $s_2 = \lambda\,e_j$
($t = \lambda$) and writing $w_j = B\,\sigma(a_j)$ gives
\begin{equation}\label{eq:2x2}
  \begin{pmatrix} 1 & 1 \\ \lambda & \lambda^e \end{pmatrix}
  \begin{pmatrix} a_j \\ w_j \end{pmatrix}
  = \begin{pmatrix} c_1 \\ c_2 \end{pmatrix}.
\end{equation}
The determinant is $\Delta = \lambda^e - \lambda$. Choosing any
$\lambda$ with $\Delta \neq 0$ (which exists for all $p > 3$) yields
a \emph{closed-form} solution:
\begin{equation}\label{eq:solution}
  a_j = \frac{\lambda^e c_1 - c_2}{\Delta}, \qquad
  w_j = \frac{-\lambda c_1 + c_2}{\Delta}.
\end{equation}

Repeating for $j = 0, \ldots, d-1$ recovers the effective key
matrix $A_{\text{eff}} := A\,P$ in $2d$ queries, where $P$ is the
publicly computable coefficient-to-stalk map at the root node.
Recovering $A$ in stalk coordinates:
\begin{equation}
  A = A_{\text{eff}} \cdot P^{-1}.
\end{equation}

With $A_{\text{eff}}$ known, $B$ is recovered by collecting $d$
additional queries with linearly independent
$z^{(i)} = \sigma(A_{\text{eff}}\,s^{(i)})$ and solving
\begin{equation}\label{eq:B-recovery}
  B = W\,Z^{-1}, \qquad
  W_i = c^{(i)} - A_{\text{eff}}\,s^{(i)}, \quad
  Z = \bigl[z^{(1)} \cdots z^{(d)}\bigr].
\end{equation}
Total query count: $2d + d = 3d$ (worst case), or $2d + 1$ when
the basis queries already provide an invertible $Z$.

\subsection{Experimental Validation}

The attack is executed against all primes
$p \in \{17, 23, 29, 47, 59, 89, 131, 251\}$, both $\sigma$
functions, and dimensions $d \in \{2, 4\}$. In every case:

\begin{itemize}
  \item $A_{\text{eff}}$ is recovered exactly (\texttt{A\_eff\_match = True}).
  \item $A$ in stalk coordinates is recovered exactly (\texttt{A\_stalk\_match = True}).
  \item $B$ is recovered exactly (\texttt{B\_match = True}).
  \item The recovered key passes 20 fresh validation queries
    (\texttt{valid = True}).
\end{itemize}

For $d = 4$ the total number of encryption queries used is $12$,
independently of $p$.

\subsection{Security Implications}
\label{sub:implications}

This result has the following consequences for the security claims
of the original paper:

\begin{enumerate}

  \item \textbf{The CPA-resistance claim is falsified for all
    parameter sets.} The attack does not require linearisation and
    is not blocked by the nonlinearity of $\sigma$. On the contrary,
    the scaling homogeneity of $\sigma$ \emph{is} the algebraic
    property that enables the attack. The statement ``linearisation
    fails because $B_v\,\sigma(A_v s_v)$ mixes key and message
    nonlinearly'' remains true but does not capture this attack
    vector.

  \item \textbf{The NL-SMIP hardness reduction is not affected.}
    The reduction from MQ establishes that recovering the key from a
    polynomial system with secret coefficients is NP-hard. The
    scaling attack does not solve NL-SMIP; it avoids it entirely by
    exploiting a structural property of $\sigma$ outside the
    polynomial formulation.

  \item \textbf{Post-quantum claims are unaffected in substance.}
    The scaling attack is purely algebraic and applies with the same
    $O(d)$ complexity classically and quantumly. It does not exploit
    a quantum speedup; it is simply faster than any exhaustive
    search.

  \item \textbf{The attack extends to all nodes simultaneously.}
    Each node's key $(A_v, B_v)$ is independent and the attack is
    node-local. Querying all $n$ nodes requires $3dn$ encryption
    queries. For a 6-node path with $d = 4$ this is 72 queries.

  \item \textbf{The attack is repairable.} Any $\sigma$ satisfying
    $\sigma(\lambda y) \neq \lambda^e \sigma(y)$ for all $e$ would
    defeat this specific attack. Extending the $\sigma$-selection
    analysis of Experiment~04 to include scaling homogeneity is the
    natural direction for a fix.

\end{enumerate}

% ================================================================
\section{Experiment 11 --- Statistical Collision Study of $g_B$}
% ================================================================

\subsection{Objective}

Obtain statistically robust estimates of the collision properties of
$g_B$ across the full parameter space, extending the spot-check of
Experiment~07 to large samples of $B$ matrices.

\subsection{Methodology}

For $d = 2$ and all primes, exact enumeration over $p^d$ inputs is
repeated for 12--24 independent draws of $B \in \GL(d,\Fp)$.
For $d = 4$, 30\,000 random inputs are sampled per $B$.
Per configuration the following are reported:
\begin{itemize}
  \item Mean image fraction $\pm$ standard deviation over the
    $B$-sample.
  \item Mean preimage multiplicity $p^d / |\Img(g_B)|$.
  \item Worst-$B$ examples (lowest image fraction).
  \item Aggregated multiplicity histogram.
\end{itemize}

\subsection{Findings and Security Implications}

The image fraction for $\sigma = \texttt{inverse}$ stabilises near
$0.63$ for $d = 2$ with very low variance ($\sigma_{\text{std}} < 0.01$
for $p \geq 59$), confirming that non-surjectivity is a stable
property of $\sigma$, not an artefact of unlucky $B$ choices.
The maximum preimage multiplicity remains bounded below 7 for
$d = 2$, $p \leq 251$, consistent with Theorem~7.1 of the original
paper. Together with Experiment~07, this experiment provides
empirical estimates of $\gamma \approx 1 - 0.63 = 0.37$ and
$L \approx 1/0.63 \approx 1.59$ for the $\sigma = \texttt{inverse}$,
$d = 2$ configuration.

% ================================================================
\section{Experiment 12 --- Key Recovery Diagnostics and Minimum Capture Analysis}
% ================================================================

\subsection{Objective}

Measure the minimum number of known-plaintext pairs required for
unique key recovery at $d = 2$ (exhaustive) and $d = 3$ (sampled),
and track how quickly false positives are eliminated as captures
accumulate.

\subsection{Methodology}

For each $\hat{A}$ candidate (exhaustive for $d = 2$, $p = 17$;
60\,000 random samples for $d = 3$), $B$ is solved from the first
$d$ captures and the resulting candidate is tested against an
increasing prefix of up to 20 captures. The \emph{first failure
index}---the smallest $k$ at which the candidate fails---is recorded.
The minimum $k$ at which exactly one candidate (the true key) remains
is denoted $k_{\min}$.

\subsection{Findings and Security Implications}

For $d = 2$, $p = 17$: the number of surviving candidates drops from
thousands at $k = d = 2$ to single digits by $k = 6$, and to exactly
1 (no false positives) by $k \approx 10$--$15$, depending on
$\sigma$. For $d = 3$ in sampled mode, unique recovery is achieved by
$k \approx 12$--$18$ captures. Each capture carries roughly
$d^2 \log_2 p$ bits of key information, consistent with the
polynomial system structure. These bounds are primarily relevant for
toy configurations; at standard parameters the exhaustive A-search
is infeasible (Experiment~09).

% ================================================================
\section{Experiment 13 --- Differential Column Attack}
% ================================================================

\subsection{Objective}

Test whether the columns of $A$ can be recovered \emph{independently}
via chosen-plaintext differential queries along coordinate axes, and
measure the resulting reduction in search complexity.

\subsection{Mathematical Basis}

For plaintexts $s = t\,e_j$, consecutive ciphertext differences are:
\begin{equation}\label{eq:diff}
  \Delta c_t = c\bigl((t+1)e_j\bigr) - c\bigl(t\,e_j\bigr)
  = a_j + B\bigl(\sigma\bigl((t+1)a_j\bigr) - \sigma(t\,a_j)\bigr).
\end{equation}
For a guessed column $\hat{a}_j$, this is a linear equation in $B$.
Collecting $d$ such differentials and denoting
$\Delta\sigma_t = \sigma((t+1)\hat{a}_j) - \sigma(t\hat{a}_j)$:
\begin{equation}
  \Delta c_t - \hat{a}_j = B\,\Delta\sigma_t.
\end{equation}
If the matrix $[\Delta\sigma_0 \cdots \Delta\sigma_{d-1}]$ is
invertible, $B$ is determined uniquely.
Each column of $A$ is recoverable by searching $p^d - 1$ nonzero
candidates for $\hat{a}_j$, reducing the total search space from
$|\GL(d,p)| \approx p^{d^2}$ to $d\,(p^d - 1)$.

\subsection{Findings}

\begin{description}
  \item[$d = 2$, $p = 17$:] $p^2 - 1 = 288$ candidates per axis
    \emph{vs.}\ $|\GL(2,17)| = 69\,360$ full matrices. Unique recovery
    in 4--6 differentials in all tested configurations.
  \item[$d = 3$, $p = 17$:] $p^3 - 1 = 4\,912$ candidates
    \emph{vs.}\ $|\GL(3,17)| \approx 1.4\times 10^7$. Unique recovery
    in 4--7 differentials.
\end{description}

\paragraph{Limitation.} If column $a_j$ has a zero in component $i$,
then $\sigma(t\,a_j)_i = 0$ for all $t$, making the $i$-th row of
$\Delta\sigma$ identically zero. The corresponding entries of $B$
remain undetermined and the attack produces false negatives for those
configurations.

\subsection{Security Implications}

The differential column attack reduces per-column search from
$p^{d^2}$ to $p^d - 1$ candidates---an exponential improvement for
$d \geq 2$. It is dominated by the $O(d)$-query attack of
Experiment~10 in practical terms, but provides an independent
recovery path useful when direct query access is restricted.

% ================================================================
\section{Discussion}
% ================================================================

\subsection{The Scaling Homogeneity as a Structural Vulnerability}

The root cause of the CPA break in Experiment~10 is the interplay
between three design choices:
\begin{enumerate}
  \item $\sigma$ is applied \emph{componentwise}, enabling
    $\sigma(\lambda y)_i = \lambda^e \sigma(y)_i$.
  \item $B$ acts \emph{linearly} on $\sigma(y)$, so
    $B\,\sigma(\lambda y) = \lambda^e\,B\,\sigma(y)$.
  \item The plaintext space includes \emph{scalar multiples} of any
    vector, allowing the attacker to query $(s, \lambda s)$ and
    decouple $y = As$ from $w = B\sigma(As)$ via a $2\times 2$
    linear system.
\end{enumerate}
Any two of these three properties in isolation would not be sufficient
for the attack. A fix must therefore break at least one of them.
The most natural option is replacing $\sigma$ with a map that lacks
componentwise scaling while preserving bijectivity and compatibility
with decryption.

\subsection{Relationship to the Original Security Proof}

The original paper's CPA resistance argument considers designs where
$\sigma$ is applied before or independently of the secret map,
showing that in those cases the key can be recovered linearly.
The final design places $\sigma$ \emph{after} the secret linear map
$A_v$, which does defeat linearisation. Experiment~10 shows that
defeating linearisation is necessary but not sufficient for CPA
security: the scaling homogeneity provides an alternative attack that
does not require the attacker to linearise the system.

The NL-SMIP hardness proof (Lemma~6.2 of the original paper)
establishes that recovering the key from a polynomial system with
secret coefficients is NP-hard. The scaling attack models a different
scenario---one in which the attacker controls the plaintexts and
exploits algebraic structure in $\sigma$. These two models are not
equivalent, and the NP-hardness of NL-SMIP does not imply CPA
security.

\subsection{Collision Properties and Decryption Correctness}

Experiments~07 and~11 together establish that $g_B$ is
non-surjective with image fraction $\approx 0.63$ for $d = 2$.
This provides empirical estimates of $\gamma \approx 0.37$ and
$L \approx 1.59$ for the correctness bound in Theorem~7.1
($\Pr[\text{ambiguity}] \leq (L-1)\,\gamma^{n-1}$). For graphs with
$n \geq 6$ nodes, decryption ambiguity is negligible in practice.

% ================================================================
\section{Conclusion}
% ================================================================

This analysis presents seven reproducible security experiments that
probe the Alaniz Cipher from multiple adversarial perspectives.
The key conclusions are:

\begin{enumerate}
  \item \textbf{Dimension $d = 1$ is trivially insecure}: the key is
    algebraically recoverable from 2 plaintext-ciphertext pairs
    (Experiment~08).

  \item \textbf{The local map $g_B$ is consistently non-surjective}
    (image fraction $\approx 0.63$) with bounded preimage
    multiplicity, validating the decryption uniqueness analysis
    (Experiments~07 and~11).

  \item \textbf{Known-plaintext exhaustive key recovery} is feasible
    for toy parameters ($d = 2$, $p \leq 29$) and requires
    approximately 10--15 captures for unique recovery; it is
    infeasible at proposed standard parameters (Experiments~09
    and~12).

  \item \textbf{A direct CPA break exists for all parameters}: the
    scaling homogeneity $\sigma(\lambda y) = \lambda^e\,\sigma(y)$
    allows recovery of the per-node key $(A_v, B_v)$ in $O(d)$
    chosen-plaintext queries, independently of $p$ (Experiment~10).
    This is the principal security finding of this analysis.

  \item \textbf{A differential column attack} reduces the per-column
    search space from $|\GL(d,p)|$ to $d\,(p^d - 1)$ candidates
    (Experiment~13).
\end{enumerate}

The scaling attack of Experiment~10 is a significant vulnerability
that requires attention before the scheme can be considered secure
under standard CPA definitions. The natural direction for a repair
is to replace the componentwise scaling-homogeneous $\sigma$ with a
map that mixes components across the $d$-dimensional fibre while
preserving the algebraic properties required for efficient decryption.

% ================================================================
\section*{Experiment Summary}
\addcontentsline{toc}{section}{Experiment Summary}
% ================================================================

\begin{table}[h]
\centering
\caption{Summary of Experiments 07--13.}
\label{tab:summary}
\renewcommand{\arraystretch}{1.35}
\begin{tabularx}{\linewidth}{c l l l X}
\toprule
Exp & Model & Attack / Analysis & Complexity & Main Result \\
\midrule
07 & ---  & $g_B$ collision characterisation
           & $O(p^d \cdot b)$
           & Image fraction $\approx 0.63$; collisions universal \\
08 & KPA  & Algebraic key recovery
           & $O(1)$ for $d{=}1$
           & $d{=}1$ trivially broken; linearisation fails at $d{\geq}2$ \\
09 & KPA  & Exhaustive A-search + linear $B$
           & $O(|\GL(d,p)|)$
           & Feasible for $p \leq 29$; intractable at standard params \\
\textbf{10} & \textbf{CPA}
           & \textbf{Scaling homogeneity attack}
           & $\bm{O(d)}$ \textbf{queries}
           & \textbf{Full key recovery at all params; CPA claim falsified} \\
11 & ---  & $g_B$ statistical multiplicity
           & $O(p^d \cdot b)$
           & Stable image fraction; worst-$B$ families identified \\
12 & KPA  & Capture count diagnostics
           & $O(|\GL(d,p)| \cdot k)$
           & Unique recovery in $\approx 10$--$15$ captures for $d{=}2$ \\
13 & CPA  & Differential column attack
           & $O(d(p^d{-}1))$
           & Search from $p^{d^2}$ reduced to $p^d{-}1$ per column \\
\bottomrule
\multicolumn{5}{l}{\small KPA = Known-Plaintext Attack; \quad CPA = Chosen-Plaintext Attack.}
\end{tabularx}
\end{table}

\bigskip
\noindent
\textit{Experiments 07--13 are implemented in
\texttt{experiments/07\_collision\_table.py} through
\texttt{experiments/13\_differential\_column\_attack.py}.
All random seeds are fixed; results are fully reproducible.}

\end{document}
